\section{Structured Sequence Transfer}
\label{sec:sequence}

\begin{figure*}[t]
  \centering
  \paperfigure{figure3_sequence_transfer}{\textwidth}
  \caption{\textbf{Independent erase/write control across repeated structured
  updates and distractor gaps.} Each E13c-R1 heat-map cell is the seed-mean
  affected-MSE gain of dual over tied control for one update-count by
  gap-length condition; annotations show the derived mean and observed seed
  range. The result concerns a shared structured event encoder, not natural
  language, a pretrained language model, an official backend, or a general
  planning agent.}
  \label{fig:sequence}
\end{figure*}

E13a is a preregistered floor/throughput calibration gate; E13c-R1 aggregates
the paired tied/dual sequence cells that followed the valid calibration path.
Across the complete update-count by distractor-gap grid, the mean affected-MSE
gain is $\CATENAResult{E13C_OVERALL_GAIN}$.  At the registered long-gap,
multi-update stress endpoint it is
$\CATENAResult{E13C_STRESS_GAIN}$ (Figure~\ref{fig:sequence}).  The complete
grid, rather than the endpoint alone, is the primary sequence bridge.

This result supports persistence of the independent erase/write advantage
behind a shared structured relational encoder.  It does not remove the oracle
conditions used by the controlled probe or establish semantic, language-model,
or agent-level transfer.
